\documentclass[11pt]{article}
\usepackage[margin=2.5cm]{geometry}
\usepackage{amsmath, amssymb, amsthm}
\usepackage{bm}
\usepackage{hyperref}
\usepackage{array}
\usepackage{booktabs}
\usepackage{multirow}

\newcommand{\R}{\mathbb{R}}
\newcommand{\SO}{\mathrm{SO}(3)}
\newcommand{\Ordn}{\mathrm{O}(3)}
\newcommand{\rmsd}{\mathrm{RMSD}}
\DeclareMathOperator*{\argmin}{arg\,min}

\title{When Are Two Conformers the Same?\\
       \large Symmetry, Chirality, and Conformer De-duplication}
\author{Notes for ChemDM}
\date{}

\begin{document}
\maketitle

\noindent
A conceptual reference for the conformer-identity question that drives the
clustering in \texttt{src/chemdm/Cluster.py}. It captures \emph{why} two
distinct-looking 3D structures can be the same chemical conformer, why plain
atom-relabelling cannot detect this, and the exact rule --- gate a reflection
on molecular achirality --- that resolves it. The running example is
nitrosocyclopentane, \texttt{O=NC1CCCC1}, whose pipeline output first raised
the question.

\section{The question}

Sample and relax the conformers of \texttt{O=NC1CCCC1}. Three survive
RMSD clustering, two of them at \emph{exactly} equal energy:

\begin{center}
\begin{tabular}{crl}
\toprule
conformer & energy (kJ/mol) & \\
\midrule
$0$ & $-58265.3736$ & \multirow{2}{*}{$\Big\}$ degenerate}\\
$1$ & $-58265.3735$ & \\
$2$ & $-58263.1315$ & distinct pucker\\
\bottomrule
\end{tabular}
\end{center}

Conformers $0$ and $1$ are \emph{not} superimposable by any rotation --- their
heavy-atom RMSD after optimal proper alignment is $0.53\,\text{\AA}$. Yet they
have identical energy. Are they two conformers or one? The answer requires
being precise about what ``the same conformer'' means.

\section{A conformer is an equivalence class}

The potential energy surface (PES) $E:\R^{3N}\to\R$ assigns an energy to every
labelled geometry $X=(x_1,\dots,x_N)$, $x_i\in\R^3$. A \emph{conformer} is a
local minimum of $E$ --- but not a single point. Physics is blind to:

\begin{itemize}
  \item \textbf{Rigid motion.} Translating or rotating the whole molecule
        changes no energy. So $X$ and $RX+t$ are the same, for $R\in\SO$,
        $t\in\R^3$.
  \item \textbf{Relabelling identical atoms.} If a permutation $\pi$ of the
        atoms maps the molecular graph onto itself (a \emph{graph
        automorphism}) --- swapping two chemically equivalent hydrogens, say
        --- then $X$ and $\pi X$ describe the same substance.
\end{itemize}

So a conformer is really an \emph{equivalence class} of geometries under the
group generated by rigid motions and the molecule's symmetry operations. The
whole question ``are $0$ and $1$ the same?'' is the question ``does the group
contain an operation taking one to the other?''. Everything below is about
which operations belong in that group.

\section{Proper vs.\ improper: the handedness invariant}

The isometries of $\R^3$ form the group $\Ordn$, which splits into two pieces
by the determinant of the linear part:
\[
  \det = +1:\quad \text{\emph{proper}} \ \ (\text{rotations, } \SO),
  \qquad
  \det = -1:\quad \text{\emph{improper}} \ \ (\text{reflections, inversions, } S_n).
\]
A proper operation preserves \emph{handedness}; an improper one reverses it.
Handedness (chirality) of a rigid point cloud is a genuine invariant: no
rotation, and no relabelling, can turn a left hand into a right hand.

This is the crux. A graph automorphism $\pi$ is a \emph{permutation} --- it
reorders labels, it does not act on space. Composing $\pi$ with a proper
rotation is still an orientation-preserving operation overall. Therefore:

\begin{center}
\emph{Relabelling plus rotation can never map a structure onto its mirror
image.}
\end{center}

If conformers $0$ and $1$ are mirror images, no amount of clever atom-matching
will superimpose them. To detect their identity you must \emph{explicitly}
include a reflection --- and whether you are \emph{allowed} to is exactly the
chirality question.

\section{The symmetry-corrected RMSD}

Fix two heavy-atom coordinate sets $A,B\in\R^{H\times 3}$ (hydrogens excluded;
they follow their heavy neighbours). Let $\mathcal{P}$ be the heavy-atom graph
automorphisms of the molecule, computed once from the 2D graph. Define
\[
  \rmsd^{\text{prop}}(A,B)
  \;=\;
  \min_{\pi\in\mathcal{P}} \ \min_{R\in\SO,\,t}
  \ \Big[\tfrac{1}{H}\textstyle\sum_i \| R a_i + t - b_{\pi(i)}\|^2\Big]^{1/2}.
\]
The inner minimisation over $(R,t)$ is solved in closed form by the
\emph{Kabsch} algorithm, using the \emph{proper} variant (a determinant-sign
correction forces $\det R=+1$, forbidding the reflected fit). This is
\texttt{kabsch\_align\_numpy} in \texttt{src/chemdm/geometry.py}, and
\texttt{\_symmetry\_corrected\_heavy\_rmsd} in \texttt{Cluster.py} is exactly
the $\min$ over $\mathcal{P}$ above.

Using the proper Kabsch is a deliberate choice: it \emph{protects chirality}.
Two genuine enantiomers of a chiral molecule must never be merged, and
$\rmsd^{\text{prop}}$ keeps them apart because it cannot cross the mirror.
The cost is that it \emph{also} keeps mirror-image conformers of an
\emph{achiral} molecule apart --- which is the over-splitting we are trying to
fix.

\section{Chirality decides whether reflection is a symmetry operation}

Here is the clean statement that resolves everything.

\begin{quote}
Reflection belongs to a molecule's symmetry group \textbf{if and only if} the
molecule is achiral.
\end{quote}

A molecule is \emph{achiral} exactly when it possesses an improper symmetry
element --- a mirror plane $\sigma$, an inversion centre $i$, or an
improper-rotation axis $S_n$ --- i.e.\ some improper isometry (possibly composed
with a relabelling) maps the molecule onto a superimposable copy of itself.
That is what ``superimposable on its mirror image'' means.

\begin{itemize}
  \item \textbf{Achiral molecule.} An improper symmetry operation
        \emph{exists}. Two geometries related by it are related by a genuine
        symmetry of the molecule, hence they are \emph{the same conformer}.
        Reflection is legitimate; include it.
  \item \textbf{Chiral molecule.} No improper symmetry operation exists.
        Reflection turns the molecule into its \emph{enantiomer} --- a
        different substance. Mirror-image conformers belong to two different
        molecules and must be kept distinct. Reflection is illegitimate;
        forbid it.
\end{itemize}

So the achiral/chiral test is not a heuristic knob. It is the precise condition
under which the reflection is a member of the equivalence group. Detect it once
from the 2D graph (no stereocentres and no stereo double bonds
$\Rightarrow$ achiral) and gate the reflected fit on it.

\section{Worked example: the nitrosocyclopentane degeneracy}

The nitroso group \mbox{O$=$N$-$C} is locally planar. The two degenerate
conformers are the cyclopentane ring puckering ``up'' on one side of that plane
and ``down'' on the other, versus the reverse --- mirror images across the
O$-$N$-$C plane. The molecule is achiral (no stereocentres), so this mirror,
composed with the ring-reversal relabelling $\text{C}_2\!\leftrightarrow\!
\text{C}_5,\ \text{C}_3\!\leftrightarrow\!\text{C}_4$, is a real $\sigma$
symmetry element.

Numerically, on the two relaxed conformers ($H=6$ heavy atoms, only two heavy
automorphisms: identity and ring-reversal):

\begin{center}
\begin{tabular}{lc}
\toprule
operation applied to align conf.\ $0$ onto conf.\ $1$ & heavy RMSD (\AA)\\
\midrule
proper rotation $+$ identity relabel        & $0.526$\\
proper rotation $+$ ring-reversal relabel   & $0.939$\\
reflection $+$ identity relabel             & $1.037$\\
\textbf{reflection $+$ ring-reversal relabel} & $\mathbf{0.000}$\\
\bottomrule
\end{tabular}
\end{center}

Neither half alone works: relabelling cannot flip handedness (rows 1--2), and a
bare reflection sends the ring atoms across the plane onto the \emph{wrong}
labels (row 3). Only the \emph{composite} improper operation --- reflection
\emph{and} the ring relabelling, i.e.\ the molecule's $\sigma$ element ---
superimposes them exactly (row 4). That composite is a symmetry operation of
this achiral molecule, which is precisely why the two are the same conformer.

\section{Are they chemically identical? Yes.}

For an isolated molecule in an achiral environment, conformers $0$ and $1$ are
\emph{the same conformer}, counted once. Being related by a molecular symmetry
operation, they are indistinguishable by every achiral observable --- energy
(exactly equal, by symmetry), dipole magnitude, moments of inertia, IR/Raman,
NMR in an achiral solvent, and all thermodynamic contributions --- and they
interconvert by a low-barrier ring flip.

The textbook analogue is \textbf{gauche-butane}. The two gauche forms
$g^{+}$ ($\phi=+60^\circ$) and $g^{-}$ ($\phi=-60^\circ$) are mirror images.
Butane is universally reported as having \emph{two} conformers --- anti and
gauche --- not three: the gauche conformer simply carries a multiplicity of
$2$. Nitrosocyclopentane is the same story: conformers $0$ and $1$ are the
$g^{+}/g^{-}$ of the ring pucker (one conformer, degeneracy $2$), and conformer
$2$ is the genuinely distinct pucker. Our pipeline already reports exactly two
conformers for butane itself, which is the sanity check.

\paragraph{The degeneracy is not lost.} Merging does not discard information:
the mirror pair contributes a \emph{symmetry factor} of $2$ to the statistical
weight (the conformer's degeneracy / symmetry number), not a second entry in
the list. If Boltzmann populations are needed, carry the multiplicity; the two
mirror wells are one state with weight $2$.

\paragraph{The one exception.} In a \emph{chiral} environment --- a chiral
solvent or catalyst, enantioselective binding, or circular-dichroism
spectroscopy --- the two conformational enantiomers respond oppositely and are
distinguishable. If the ensemble feeds such a calculation, keep both. For every
achiral-observable purpose (energies, populations, ordinary spectra, matching
to CREST references) they are redundant and one should be removed.

\section{The de-duplication rule, precisely}

\begin{enumerate}
  \item Compute the heavy-atom graph automorphisms $\mathcal{P}$ once from the
        molecule (identity included).
  \item Determine achirality once: no assigned or potential stereocentres and
        no stereo double bonds.
  \item Two relaxed conformers $A,B$ are the same iff, after energy
        pre-filtering,
        \[
          \min\Big(\rmsd^{\text{prop}}(A,B),\
          \underbrace{[\,\text{achiral}\,]\cdot \rmsd^{\text{refl}}(A,B)}_{\text{only if achiral}}\Big)
          \ \le\ \tau,
        \]
        where $\rmsd^{\text{refl}}$ is the same minimisation but allowing the
        reflected Kabsch fit, and $[\,\text{achiral}\,]$ is $1$ for an achiral
        molecule and $+\infty$ (disabled) for a chiral one.
\end{enumerate}

For a chiral molecule the second term is switched off and the behaviour is
identical to today's proper-only clustering: genuine enantiomers are never
merged. For an achiral molecule the reflected fit is enabled, and it can only
ever collapse pairs that a real improper symmetry operation of the molecule
connects --- exactly the conformational enantiomers, and nothing else. (In the
data above, the distinct pucker, conformer $2$, stays $0.33\,\text{\AA}$ from
the mirror pair even \emph{with} reflection: it is not over-merged.)

\paragraph{One subtlety worth a footnote.} ``Achiral molecule'' refers to the
\emph{configuration} (the 2D graph): no fixed stereocentres. An individual
puckered \emph{conformation} may still be geometrically chiral (conf.\ $0$ has
$C_1$ point-group symmetry). That is fine and expected --- it is the whole
reason the two conformers look different. The molecule being configurationally
achiral is what guarantees the two chiral conformations are mirror-related and
interconvert, hence the same species.

\end{document}
